Mineral scale that narrows the production choke of an oil well grows over hours to weeks, yet it is
usually noticed late from trend charts. Machine-learning studies on the public Petrobras 3W dataset detect
such events once they have begun, mostly with the same wells in training and test sets, and none forecasts
when the restriction becomes critical. This paper formulates choke-scale forecasting as discrete-time
survival analysis: from each 6-h window of 1-min production-sensor data, a model estimates the probability
that a self-referenced restriction index rises by 5\,\% within horizons of 1 to 72~h, with censored
instances handled explicitly. The proposed model, PI-PatchSurv, fuses a channel-independent patch
Transformer with a physics-feature branch through a learned gate and outputs a monotone risk curve; it is
trained with a censoring-aware likelihood and a monotonicity penalty, and its uncertainty is quantified with
a five-member ensemble and split-conformal intervals. Twelve models were compared on 3W version 2.0.0 by
leave-one-well-out validation (24 real events in six wells), with alarm thresholds fixed on separate
calibration wells. Recurrent forecasters ranked events best, with an area under the precision--recall
curve at 24~h of up to 0.74 against 0.53 for a physics threshold rule, and caught up to 17 of 24 events a
median of 2.9~h ahead. PI-PatchSurv ranked lower (0.64) but gave the best-calibrated intervals (coverage
0.87 at the 90\,\% level against 0.33--0.54), the smallest time-to-event error (6.8~h) and no loss when
three of nine sensors failed. Forecasting the time to choke restriction on unseen wells is therefore
feasible, with lead times bounded by the short recordings.
